\section{Ex-post shielding: prices versus incomes}
\label{sec:shield}

The crisis instruments of the recent literature are a retail price cap and an
income transfer. Both stabilise the crisis, and
Table~\ref{tab:summary} shows they do so comparably. They differ on the margin
this paper adds.

\input{output/tab_summary_cpi_import.tex}

\paragraph{Cap versus transfer: the non-neutrality.}\label{sec:cap}\label{sec:transfer}

Under the cap the peak green-share response falls from $8.0$ pp to essentially
zero ($\approx-0.05$ pp) under \emph{both} shocks: the margin is not attenuated; it is
switched off. The mechanism is direct: the cap holds $P^{E}_{B}$ at its
pre-crisis level, and $P^{E}_{B}$ determines the return to adoption. The full fiscal responses under both shocks are in
Figures~\ref{fig:policy_price} and~\ref{fig:policy_supply} (no policy, price cap,
Slutsky transfer).

\begin{figure}[H]\centering
\figtitle{Price shock: fiscal responses}
\includegraphics[width=0.92\textwidth]{output/fig2_policy_price.png}
\label{fig:policy_price}
\fignotes{Adoption open. No policy (solid), price cap (dashed), Slutsky transfer
(dash-dotted).}
\end{figure}

\begin{figure}[H]\centering
\figtitle{Supply shock: fiscal responses}
\includegraphics[width=0.92\textwidth]{output/fig2_policy_supply.png}
\label{fig:policy_supply}
\fignotes{Adoption open. No policy (solid), price cap (dashed), Slutsky transfer
(dash-dotted).}
\end{figure}

The Slutsky transfer leaves the green-share response essentially unchanged
($8.5$ pp against $8.0$ with no policy) while delivering more output
stabilisation than the cap at similar fiscal cost ($\sum y=-25.9$ vs.\ $-53.1$;
gross fiscal $49.2$ vs.\ $54.1$). This preservation follows directly from the transfer's
lump-sum indexing (Section~\ref{sec:bop}): because it does not condition on the
household's current durable, switching forfeits nothing and the relative-price
incentive to adopt is left intact.

The sixth panel reports fiscal outlays $X$. Under the price shock the cap's outlay
stays positive; under the supply shock it turns negative in the recovery
(Figure~\ref{fig:policy_supply}). With the retail price frozen, household energy
demand cannot adjust, so the fixed-quantity market clears through $P^{E*}$ alone: the
clearing price overshoots during the shock and falls below its pre-crisis level
afterwards. The cap pays the gap between $P^{E*}$ and the frozen retail price, so once
$P^{E*}$ drops below the freeze the scheme collects rather than pays. This reversal is
larger than an integrated market's demand elasticity would produce, so the
supply-shock cap magnitudes should be read with caution.

\paragraph{Two margins the frozen price governs.}\label{sec:exp_switching}

Adoption is not the only margin the frozen relative price governs. The
open-economy identity $y=c_H+c_H^{*}$ splits home-good output into domestic demand
for the home good and exports (Figure~\ref{fig:output_decomp}). Under no policy the energy shock lifts exports: the
relative price of the home good falls and foreign demand rises, a partial offset to
the domestic-demand collapse (cumulatively $+7.4$ against a domestic $-72.3$). The
cap switches that offset off almost entirely ($+0.7$): holding the retail energy
price fixed suppresses the relative-price movement expenditure switching needs,
exactly as it suppresses the one adoption needs. The transfer leaves the relative
price free and preserves the export channel nearly intact ($+6.1$) while supporting
domestic demand ($-32.0$). The cap and the transfer stabilise comparable amounts of
output, but the cap does so by suppressing an adjustment margin and the transfer by
supporting income: the same distinction the paper draws on the adoption share, read
now through the trade balance.

\begin{figure}[H]\centering
\figtitle{Output response: domestic demand and exports}
\IfFileExists{output/fig_output_domestic_external_import.pdf}{\includegraphics[width=\textwidth]{output/fig_output_domestic_external_import.pdf}}{\textit{[figure output/fig\_output\_domestic\_external\_import.pdf: run run\_output\_decomposition.py]}}
\label{fig:output_decomp}
\fignotes{Brown-price shock, by policy. Output decomposed as $y=c_H+c_H^{*}$,
domestic demand for the home good (blue) plus exports (green). Level deviations
$\times100$.}
\end{figure}

\paragraph{Partial caps.}\label{sec:dial}

Table~\ref{tab:summary} uses $\tau^{E}=1$, a complete price freeze;
Table~\ref{tab:dose} sweeps $\tau^{E}\in[0,1]$. Under the \emph{price} shock the
adoption loss is almost exactly linear in $\tau^{E}$. Under the \emph{supply}
shock it is convex: small caps barely touch adoption and the collapse is
concentrated near the full freeze. A partial shield is therefore much less damaging to the transition when supply is
inelastic.

At the full freeze under the supply shock the fiscal cost is negative
(Table~\ref{tab:dose}), for the reason given in Section~\ref{sec:cap}: in the recovery
the clearing price falls below the frozen retail level, so the cap collects rather
than pays.

\input{output/tab_dose_cpi_import.tex}

\begin{figure}[H]\centering
\figtitle{Price cap: response across intensity}
\includegraphics[width=\textwidth]{output/fig6_dose_response.png}
\label{fig:dose}
\fignotes{Cap intensity $\tau^E$ (red) against the Slutsky transfer benchmark
(blue dashed), for both shocks. Rows: price shock (top), supply shock (bottom).}
\end{figure}

\paragraph{The green subsidy.}\label{sec:greensub}

The green subsidy is a crisis instrument in its own
right (Figure~\ref{fig:green_subsidy}): the government pays a fraction of the
switching cost during the acute crisis. It roughly doubles the adoption response
(peak $D^{G}$ of $18.9$ vs.\ $8.0$ pp; switchers $2.6$ vs.\ $1.0$ pp) at trivial
fiscal cost ($2.9$ against the cap's $54.1$), yet its CEV is the worst in
Table~\ref{tab:summary_price_import}, $-1.77\%$, below laissez-faire, and it
deepens the output loss ($\sum y=-77.7$). Paying households to incur the real
switching cost $\psi_g D^{sw}$, booked as an import (Section~\ref{sec:bop}), in
the middle of a consumption crisis destroys welfare even as it greens the economy.

\begin{figure}[H]\centering
\figtitle{The green subsidy}
\IfFileExists{output/fig_green_subsidy_import.pdf}{\includegraphics[width=\textwidth]{output/fig_green_subsidy_import.pdf}}{\textit{[figure: run run\_green\_subsidy.py]}}
\label{fig:green_subsidy}
\fignotes{Brown-price shock, against laissez-faire, the price cap and the Slutsky
transfer. The subsidy pays a fraction of the switching cost $\psi_g$ during the
acute crisis.}
\end{figure}

The subsidy here is a crisis instrument, paid from the deficit; its use as a recycling channel for permanent carbon revenue is a different exercise, taken up in Appendix~\ref{sec:recycle}.

\paragraph{Who to shield: targeting and incidence.}\label{sec:flat}

Without policy the impatient (high-MPC) households lose far more than the patient
(Table~\ref{tab:cev}). This gap makes the incidence of the transfer a first-order
design question. The Slutsky transfer indexes the lump sum to each household's
pre-crisis brown-energy use. Replacing it with an \emph{untargeted} flat transfer
of the same fiscal cost ($49.2$ in both, over $H=24$) cuts the welfare
loss by about a third, from a CEV of $-0.88\%$ to $-0.58\%$ (Table~\ref{tab:summary_price_import}).
It even improves output ($\sum y=-22.2$ vs.\ $-25.9$), while leaving adoption
intact ($8.5$ pp). The gain is entirely distributional.

The mechanism is that energy-indexed targeting is regressive \emph{in this model}.
With homothetic CES energy demand the energy share is constant across the wealth
distribution, so brown-energy use rises with total consumption, hence with
permanent income and patience. Indexing the transfer to pre-crisis energy use
therefore hands the largest cheques to the low-MPC, wealthy households, where the
insurance value per euro is lowest, while a flat transfer reaches the high-MPC
households the crisis hits hardest. The result speaks to the 2022--23 European energy
``brakes,'' which were indexed to historical consumption: within the model the
indexing is the source of the inefficiency rather than a refinement of it.

The result depends on the homothetic energy nest: if
energy were a necessity with a declining Engel curve, the empirically standard
case, low-income households would spend a \emph{larger} share and the sign could
reverse. The statement is conditional: \emph{under constant energy shares,
energy-indexed targeting is dominated by an equal-value flat transfer}.

\input{output/tab_cev_cpi_import.tex}
